% ---------------------------------------------------------------------------
% Kuramoto oscillator environments
% ---------------------------------------------------------------------------
\begingroup
% Extra interword stretch for paragraphs with long, unbreakable identifiers
% (local to this chapter).
\setlength{\emergencystretch}{2em}
\chapter{Kuramoto Oscillator Networks}\label{chap:kuramoto}

The Kuramoto environments simulate a network of coupled phase oscillators that
a controller has to synchronise. They are used to study learning-based control
of network synchronisation: shaping the coupling of a network (coupling
design), steering oscillators with bounded inputs, tracking a common frequency,
and comparing learned controllers with classical results such as the
synchronisation transition of the Kuramoto model. Two interchangeable backends
implement the same model:

\begin{itemize}
  \item \code{KuramotoOscillatorEnv} (module \code{env\_lib.kos\_env.kuramoto\_env}):
        NumPy, one network per environment;
  \item \code{KuramotoOscillatorEnvTorch} (module
        \code{env\_lib.kos\_env.kuramoto\_env\_torch}): PyTorch, \code{n\_agents}
        independent networks simulated in parallel on the CPU or a GPU; requires
        the \code{torch} extra.
\end{itemize}

With identical initial states and actions both backends produce the same
trajectories up to \code{float32} rounding.

% ---------------------------------------------------------------------------
\section{Model}\label{sec:kuramoto-model}

The state consists of the phases $\theta_1, \dots, \theta_N$ of $N$ oscillators
with natural frequencies $\omega_i$. The controlled dynamics are
\begin{equation}
  \dot\theta_i = \omega_i + a_i + c \sum_{j=1}^{N} K_{ij} \sin(\theta_j - \theta_i),
  \qquad i = 1, \dots, N,
  \label{eq:kuramoto}
\end{equation}
where $a_i$ is the control input of oscillator $i$, $K$ the coupling matrix and
$c = 1$, or $c = 1/N$ with \code{normalize\_coupling=True} (the classical
$K/N$ scaling). Positive $K_{ij}$ are attractive and negative $K_{ij}$
repulsive couplings; the diagonal of $K$ has no effect.

\paragraph{Integration.} One environment step advances the system by
$\Delta t$ = \code{dt} with the explicit Euler method,
$\theta \leftarrow \theta + \Delta t\, f(\theta)$, or with the classical
fourth-order Runge--Kutta method (\code{integration\_method="rk4"}),
\begin{gather*}
  \theta \leftarrow \theta + \tfrac{\Delta t}{6}\,(k_1 + 2k_2 + 2k_3 + k_4), \\
  k_1 = f(\theta),\quad
  k_2 = f(\theta + \tfrac{\Delta t}{2} k_1),\quad
  k_3 = f(\theta + \tfrac{\Delta t}{2} k_2),\quad
  k_4 = f(\theta + \Delta t\, k_3),
\end{gather*}
where $f$ is the right-hand side of \eqref{eq:kuramoto} with the action held
constant during the step. If \code{noise\_std} $= \sigma > 0$, independent
Gaussian noise $\mathcal{N}(0, \sigma^2)$ is added to every phase after the
integration step (the noise is not scaled by $\sqrt{\Delta t}$). Finally the
phases are wrapped to $[-\pi, \pi)$.

\paragraph{Synchronisation measures.} The complex order parameter
\begin{equation}
  r\,e^{i\psi} = \frac{1}{N} \sum_{j=1}^{N} e^{i\theta_j}
  \label{eq:order-parameter}
\end{equation}
has magnitude $r \in [0, 1]$: $r = 1$ for identical phases and $r \approx 0$
for incoherent phases; $\psi$ is the mean phase. The phase coherence is
\begin{equation}
  C = \exp\bigl(-\operatorname{Var}(\theta)\bigr),
  \label{eq:coherence}
\end{equation}
where $\operatorname{Var}$ is the population variance (normalised by $N$) of
the phases wrapped to $[-\pi, \pi)$. For frequency tracking the environment
uses the mean absolute frequency error
\begin{equation}
  E_f = \frac{1}{N} \sum_{i=1}^{N} \bigl|\dot\theta_i - \Omega\bigr|,
  \label{eq:freq-error}
\end{equation}
with target frequency $\Omega$ = \code{target\_frequency}; $\dot\theta_i$ is
evaluated at the start of the step, with the new action applied (for RK4 this
is $k_1$).

% ---------------------------------------------------------------------------
\section{Coupling modes and topologies}\label{sec:kuramoto-coupling}

The network topology is an undirected graph with adjacency matrix $A$. Only
connected pairs are coupled. The available topologies are listed in
Table~\ref{tab:kuramoto-topologies}; \code{adj\_matrix} overrides
\code{topology}, which is then reported as \code{"custom"}.

\begin{table}[htbp]
  \centering\small
  \caption{Network topologies. $M$ is the number of undirected edges.}
  \label{tab:kuramoto-topologies}
  \begin{tabularx}{\textwidth}{@{}lXl@{}}
    \toprule
    \code{topology} & Edges & $M$ \\
    \midrule
    \code{"fully\_connected"} & every pair of oscillators & $N(N-1)/2$ \\
    \code{"ring"} & $i$ is connected to $i \pm 1 \bmod N$ & $N$ (for $N \ge 3$) \\
    \code{"star"} & oscillator 0 is connected to all others & $N - 1$ \\
    \code{"random"} & every ordered pair is drawn with probability $1/2$ and the
      matrix is symmetrised with a logical OR (a pair is connected with
      probability $3/4$); the draw uses a local generator seeded with
      \code{topology\_seed} & random \\
    \code{"custom"} & non-zero entries of \code{adj\_matrix} (either direction) & varies \\
    \bottomrule
  \end{tabularx}
\end{table}

\paragraph{Dynamic coupling} (\code{coupling\_mode="dynamic"}, default). The
coupling strength of every edge is part of the action. The edges are listed in
row-major order of the upper triangle of $A$; the list is available as
\code{env.coupling\_indices} (pairs $(i, j)$ with $i < j$) and its length as
\code{env.n\_couplings} $= M$. The strength $k_e$ of edge $e = (i, j)$ sets
$K_{ij} = K_{ji} = k_e$. A non-symmetric \code{adj\_matrix} triggers a warning,
because every connected pair is one undirected edge with a single strength.

\paragraph{Constant coupling} (\code{coupling\_mode="constant"}). $K$ is fixed:
it is \code{constant\_coupling\_matrix} when given (the diagonal is set to zero)
and \code{coupling\_strength}~$\cdot A$ otherwise. Only the control inputs are
actions. A \code{constant\_coupling\_matrix} passed in dynamic mode is ignored
with a warning.

% ---------------------------------------------------------------------------
\section{Spaces}\label{sec:kuramoto-spaces}

\paragraph{Action.} A \code{float32} \code{Box} with the layout of
Table~\ref{tab:kuramoto-action}. Actions are clipped to the bounds; an action
of the wrong shape or with non-finite values raises \code{ValueError}.

\begin{table}[htbp]
  \centering\small
  \caption{Action layout.}
  \label{tab:kuramoto-action}
  \begin{tabularx}{\textwidth}{@{}llX@{}}
    \toprule
    Indices & Content & Bounds \\
    \midrule
    \code{0 : N} & control inputs $a_1, \dots, a_N$ & \code{control\_input\_range}, default $[-1, 1]$ \\
    \code{N : N+M} & coupling strengths $k_1, \dots, k_M$ (dynamic mode only) &
      \code{coupling\_range}, default $[0, 5]$ \\
    \bottomrule
  \end{tabularx}
\end{table}

\paragraph{Observation.} An unbounded \code{float32} \code{Box} of size $3N + M$
in dynamic mode and $3N$ in constant mode (Table~\ref{tab:kuramoto-obs}). The
default \envid{KuramotoOscillator-v0} ($N = 10$, fully connected, $M = 45$) has
an action of size 55 and an observation of size 75.

\begin{table}[htbp]
  \centering\small
  \caption{Observation layout (dynamic mode; in constant mode the coupling block is absent and the control inputs occupy \code{2N : 3N}).}
  \label{tab:kuramoto-obs}
  \begin{tabularx}{\textwidth}{@{}lX@{}}
    \toprule
    Indices & Content \\
    \midrule
    \code{0 : N} & phases $\theta_i$, wrapped to $[-\pi, \pi)$ \\
    \code{N : 2N} & natural frequencies $\omega_i$ \\
    \code{2N : 2N+M} & coupling strengths of the last action (sampled at reset) \\
    \code{2N+M : 3N+M} & control inputs of the last action (zero after reset) \\
    \bottomrule
  \end{tabularx}
\end{table}

% ---------------------------------------------------------------------------
\section{Reward and episode end}\label{sec:kuramoto-reward}

The reward is computed after the step from the new phases. The reward type is
selected with \code{reward\_type}:
\begin{align*}
  \text{\code{"order\_parameter"}:}\quad & R = r, \\
  \text{\code{"phase\_coherence"}:}\quad & R = C, \\
  \text{\code{"combined"}:}\quad & R = r + C, \\
  \text{\code{"frequency\_synchronization"}:}\quad & R = -E_f,
\end{align*}
with $r$, $C$ and $E_f$ from \eqref{eq:order-parameter}--\eqref{eq:freq-error}.
When $r$ exceeds \code{sync\_threshold} (default 0.99) the episode terminates
(\code{terminated=True}) and, for all reward types except
\code{"frequency\_synchronization"}, \code{sync\_bonus} (default 10) is added
to the reward. A threshold above 1 disables termination on synchronisation.
\code{truncated} becomes \code{True} after \code{max\_steps} steps.

% ---------------------------------------------------------------------------
\section{Parameters}\label{sec:kuramoto-params}

Both backends accept the parameters of Table~\ref{tab:kuramoto-params}; all
have defaults. Invalid values raise \code{ValueError} or \code{TypeError}.

\begingroup\small
\begin{longtable}{@{}>{\raggedright\arraybackslash}p{0.29\textwidth}>{\raggedright\arraybackslash}p{0.21\textwidth}>{\raggedright\arraybackslash}p{0.44\textwidth}@{}}
  \caption{Constructor parameters of \code{KuramotoOscillatorEnv} and \code{KuramotoOscillatorEnvTorch}.}\label{tab:kuramoto-params}\\
  \toprule
  Parameter & Default & Description \\
  \midrule
  \endfirsthead
  \toprule
  Parameter & Default & Description \\
  \midrule
  \endhead
  \bottomrule
  \endfoot
  \code{n\_oscillators} & \code{10} & Number of oscillators $N$ ($\ge 2$). \\
  \code{n\_agents} & \code{1} & PyTorch: number of systems simulated in parallel. NumPy: only sets the length of \code{info["agent\_rewards"]}. \\
  \code{dt} & \code{0.01} & Integration time step ($> 0$). \\
  \code{max\_steps} & \code{1000} & Episode length; reaching it sets \code{truncated}. \\
  \code{coupling\_range} & \code{(0.0, 5.0)} & Bounds of the coupling actions and of the random initial coupling strengths. \\
  \code{control\_input\_range} & \code{(-1.0, 1.0)} & Bounds of the control inputs $a_i$. \\
  \code{natural\_freq\_range} & \code{(0.5, 2.0)} & Natural frequencies are drawn uniformly from this range at every reset. \\
  \code{device} & \code{"cpu"} & PyTorch: \code{"cpu"}, \code{"cuda"}, \code{"cuda:k"}, a \code{torch.device}, or \code{"auto"} (CUDA when available). NumPy: ignored. \\
  \code{render\_mode} & \code{None} & \code{None}, \code{"rgb\_array"} or \code{"human"}. \\
  \code{integration\_method} & \code{"euler"} & \code{"euler"} or \code{"rk4"}. \\
  \code{reward\_type} & \code{"order\_parameter"} & \code{"order\_parameter"}, \code{"phase\_coherence"}, \code{"combined"} or \code{"frequency\_synchronization"}. \\
  \code{noise\_std} & \code{0.0} & Standard deviation of the phase noise added after each step. \\
  \code{topology} & \code{"fully\_connected"} & \code{"fully\_connected"}, \code{"ring"}, \code{"star"}, \code{"random"} or \code{"custom"}. \\
  \code{adj\_matrix} & \code{None} & Custom $(N, N)$ adjacency matrix (array or tensor); overrides \code{topology}. \\
  \code{target\_frequency} & \code{1.0} & Target $\Omega$ of the frequency-synchronisation reward. \\
  \code{coupling\_mode} & \code{"dynamic"} & \code{"dynamic"} (couplings are actions) or \code{"constant"}. \\
  \code{constant\_coupling\_matrix} & \code{None} & Fixed $(N, N)$ coupling matrix of constant mode. \\
  \code{coupling\_strength} & \code{1.0} & Uniform strength of constant mode without an explicit matrix ($K = $ \code{coupling\_strength}$\,\cdot A$). \\
  \code{normalize\_coupling} & \code{False} & Divide the coupling sum by $N$. \\
  \code{topology\_seed} & \code{42} & Seed of the local generator that builds the \code{"random"} topology. \\
  \code{sync\_threshold} & \code{0.99} & Order parameter above which the episode terminates. \\
  \code{sync\_bonus} & \code{10.0} & Reward bonus on synchronisation. \\
  \code{render\_agent} & \code{0} & PyTorch only: index of the system shown by \code{render()}. \\
\end{longtable}
\endgroup

\paragraph{Initial state and reset options.} \code{reset()} draws phases
uniformly from $[-\pi, \pi)$, natural frequencies uniformly from
\code{natural\_freq\_range} and, in dynamic mode, coupling strengths uniformly
from \code{coupling\_range}; the control inputs start at zero. Any part can be
fixed with \code{reset(options=...)} using the keys \code{"phases"},
\code{"natural\_frequencies"} (shape $(N,)$) and \code{"coupling\_strengths"}
(shape $(M,)$, dynamic mode only). The PyTorch backend also accepts arrays of
shape \code{(n\_agents, size)} with one row per system.

% ---------------------------------------------------------------------------
\section{Info dictionary}\label{sec:kuramoto-info}

\begin{table}[htbp]
  \centering\small
  \caption{Keys of the info dictionary. The PyTorch backend returns per-system
    NumPy arrays with a leading dimension \code{n\_agents} for the quantities
    marked with an asterisk.}
  \label{tab:kuramoto-info}
  \begin{tabularx}{\textwidth}{@{}lX@{}}
    \toprule
    Key & Content \\
    \midrule
    \code{order\_parameter}* & order parameter $r$ after the step \\
    \code{phase\_coherence}* & phase coherence $C$ after the step \\
    \code{phases}* & wrapped phases, shape $(N,)$ \\
    \code{natural\_frequencies}* & natural frequencies $\omega_i$ \\
    \code{coupling\_matrix}* & coupling matrix $K$ used in the step, shape $(N, N)$ \\
    \code{dphases\_dt}* & phase velocities $\dot\theta$ at the start of the step (not in the reset info) \\
    \code{step\_count} & number of steps taken in the episode \\
    \code{device} & device of the simulation (\code{"cpu"} for NumPy) \\
    \code{n\_agents} & value of \code{n\_agents} \\
    \code{agent\_rewards} & \code{float64} array of shape \code{(n\_agents,)}; NumPy: copies of the reward; PyTorch: the reward of every system, including its bonus (not in the reset info) \\
    \bottomrule
  \end{tabularx}
\end{table}

% ---------------------------------------------------------------------------
\section{PyTorch backend}\label{sec:kuramoto-torch}

\code{KuramotoOscillatorEnvTorch} simulates \code{n\_agents} independent
networks with the same topology in one batch of tensors. All state tensors have
a leading batch dimension.

\begin{itemize}
  \item \code{step(action)} accepts an action of shape \code{(act\_dim,)},
        applied to every system, or \code{(n\_agents, act\_dim)} with one row
        per system, as a NumPy array or a tensor.
  \item The returned observation, scalar reward and \code{terminated} flag
        refer to system 0; \code{truncated} is shared. The per-system values are
        in \code{info} (Table~\ref{tab:kuramoto-info}).
  \item \code{get\_batch\_observations()} returns the observations of all
        systems as a tensor of shape \code{(n\_agents, obs\_dim)} on
        \code{env.device}.
  \item \code{get\_batch\_rewards(dphases\_dt=None, include\_bonus=False)}
        returns the per-system rewards of the current state as a tensor of shape
        \code{(n\_agents,)}; the synchronisation bonus is only included with
        \code{include\_bonus=True} (as in \code{info["agent\_rewards"]}).
  \item Systems are not reset individually: a system that synchronises keeps
        evolving (and keeps earning its bonus) until the whole batch is reset.
\end{itemize}

Computations run in \code{float32}. For large batches or networks select a GPU
with \code{device="cuda"} or \code{device="auto"}.

% ---------------------------------------------------------------------------
\section{Rendering}\label{sec:kuramoto-render}

Figure~\ref{fig:kuramoto} shows the dashboard ($1000 \times 560$ pixels). The
phase portrait on the left places every oscillator on the unit circle at angle
$\theta_i$, coloured by its natural frequency (colour bar), with a glow and a
fading trail of its recent phases. Couplings are drawn as chords whose width
and opacity scale with $|K_{ij}|$; repulsive couplings use the theme's negative
colour. The arrow from the origin is the order parameter $r e^{i\psi}$. On the
right, the upper panel shows $r(t)$ and the phase coherence over the episode
with the synchronisation threshold as a dashed line, and the lower panel a
raster of $\sin\theta_i$ over the most recent steps. The header reports the
step, time, $r$ and the reward. The PyTorch backend shows system
\code{render\_agent}.

\begin{figure}[htbp]
  \centering
  \includegraphics[width=\figwidth]{kuramoto.png}
  \caption{Kuramoto dashboard (16 oscillators on a ring, dynamic coupling,
    light theme). Left: phase portrait with coupling chords, trails and the
    order-parameter vector. Right: order parameter and phase coherence over
    time, and the phase raster $\sin\theta_i$.}
  \label{fig:kuramoto}
\end{figure}

% ---------------------------------------------------------------------------
\section{Registered ids}\label{sec:kuramoto-ids}

\begin{table}[htbp]
  \centering\small
  \caption{Kuramoto ids (all other parameters keep their defaults).}
  \label{tab:kuramoto-ids}
  \begin{tabularx}{\textwidth}{@{}>{\raggedright\arraybackslash}p{0.52\textwidth}X@{}}
    \toprule
    Id & Configuration \\
    \midrule
    \envid{KuramotoOscillator-v0} & NumPy, $N = 10$, dynamic coupling \\
    \envid{KuramotoOscillator-v1} & NumPy, $N = 5$, dynamic coupling \\
    \envid{KuramotoOscillator-Constant-v0} & NumPy, $N = 6$, constant coupling \\
    \envid{KuramotoOscillator-FreqSync-Constant-v0} & NumPy, $N = 6$, constant coupling, frequency reward, $\Omega = 2$ \\
    \envid{KuramotoOscillatorTorch-v0} & PyTorch, $N = 10$, one system \\
    \envid{KuramotoOscillatorTorch-v1} & PyTorch, $N = 8$, 4 systems, CPU \\
    \envid{KuramotoOscillatorTorch-v2} & PyTorch, $N = 10$, RK4, \code{device="auto"} \\
    \envid{KuramotoOscillatorTorch-Constant-v0} & PyTorch, $N = 6$, constant coupling \\
    \envid{KuramotoOscillatorTorch-FreqSync-Constant-v0} & PyTorch, $N = 6$, constant coupling, frequency reward, $\Omega = 2$ \\
    \bottomrule
  \end{tabularx}
\end{table}

% ---------------------------------------------------------------------------
\section{Usage}\label{sec:kuramoto-usage}

The first example drives a network with constant coupling to synchronisation
with a proportional phase-alignment input $a_i = r \sin(\psi - \theta_i)$
computed from the observation:

\begin{pycode}
import numpy as np
import env_lib

env = env_lib.KuramotoOscillatorEnv(
    n_oscillators=12, coupling_mode="constant", coupling_strength=0.5, max_steps=1000
)
obs, info = env.reset(seed=0)
n = env.n_oscillators
while True:
    phases = obs[:n].astype(np.float64)
    z = np.exp(1j * phases).mean()                      # r exp(i psi)
    action = np.abs(z) * np.sin(np.angle(z) - phases)   # phase alignment
    obs, reward, terminated, truncated, info = env.step(action.astype(np.float32))
    if terminated or truncated:
        break
print(info["step_count"], info["order_parameter"])      # synchronised after 51 steps
env.close()
\end{pycode}

The second example runs four systems in parallel on the PyTorch backend, one
coupling level per system:

\begin{pycode}
import numpy as np
import env_lib

env = env_lib.make("KuramotoOscillatorTorch-v1")   # 8 oscillators, 4 systems
core = env.unwrapped
obs, info = env.reset(seed=0)
print(core.get_batch_observations().shape)         # torch.Size([4, 52])

n, act_dim = core.n_oscillators, core.action_space.shape[0]
actions = np.zeros((core.n_agents, act_dim), dtype=np.float32)
actions[:, n:] = np.linspace(1.0, 4.0, core.n_agents)[:, None]
for _ in range(300):
    obs, reward, terminated, truncated, info = env.step(actions)
print(info["order_parameter"])                     # one value per system
print(info["agent_rewards"])                       # per-system rewards incl. bonus
env.close()
\end{pycode}

The example script \code{examples/kuramoto\_demo.py} ramps the coupling
strengths and adds phase alignment, and records the transition from incoherence
to synchronisation:

\begin{shellcode}
python examples/kuramoto_demo.py --save renders/kuramoto.gif
python examples/kuramoto_demo.py --backend torch --n-agents 4 --theme light
python examples/kuramoto_demo.py --mode constant --topology two_clusters --render-mode human
\end{shellcode}

% ---------------------------------------------------------------------------
\section{Reproducibility and performance}\label{sec:kuramoto-perf}

\paragraph{Seeding.} \code{reset(seed=s)} seeds \code{self.np\_random}, from
which the NumPy backend draws the initial state and the phase noise. The
PyTorch backend seeds its own \code{torch.Generator} from
\code{self.np\_random} at every reset; neither backend touches the global
random state. The random topology is built once from \code{topology\_seed}. The
two backends draw from different generators, so the same seed gives different
initial states; to compare them, pass the initial state through
\code{reset(options=...)}.

\paragraph{Performance.} Table~\ref{tab:kuramoto-perf} lists step times
measured with \code{benchmarks/benchmark\_envs.py} on the development machine
(random actions, no rendering). The coupling term is evaluated with two
matrix--vector products instead of $N^2$ sine evaluations. Rendering in
\code{"rgb\_array"} mode takes about 25--35\,ms per frame and dominates the
cost of recorded rollouts; use \code{render\_every} in \code{record\_episode} to
thin out long episodes. Use Euler integration for speed and RK4 when accuracy
at larger \code{dt} matters, and a GPU for large batches.

\begin{table}[htbp]
  \centering\small
  \caption{Step time in milliseconds on the development machine.}
  \label{tab:kuramoto-perf}
  \begin{tabular}{@{}lrr@{}}
    \toprule
    Case & Before 1.0 & Version 1.0 \\
    \midrule
    NumPy, $N = 10$ & 0.17 & 0.04 \\
    NumPy, $N = 50$ & 2.36 & 0.07--0.11 \\
    NumPy, RK4, $N = 50$ & 7.1 & 0.1 \\
    PyTorch (CPU), $N = 50$, 8 systems & 10.3 & 0.3--0.4 \\
    \bottomrule
  \end{tabular}
\end{table}

% ---------------------------------------------------------------------------
\section{Changes in 1.0}\label{sec:kuramoto-changes}

The behaviour of the Kuramoto environments changed in several ways; see
Appendix~\ref{app:migration} for migration advice.

\begin{itemize}
  \item The PyTorch backend had a sign error that made the coupling repulsive
        ($\sin(\theta_i - \theta_j)$ instead of $\sin(\theta_j - \theta_i)$).
        Both backends now implement \eqref{eq:kuramoto} and produce the same
        trajectories; normalisation by $N$ is opt-in for both.
  \item The coupling uses the full matrix; the former NumPy loop skipped entries
        $K_{ij} \le 0$. Default settings are unaffected.
  \item Synchronisation sets \code{terminated} and the step limit sets
        \code{truncated} (previously both were reported as \code{terminated});
        \code{sync\_threshold} and \code{sync\_bonus} are configurable.
  \item Constant mode without a matrix uses
        \code{coupling\_strength}$\,\cdot A$, so the \code{*-Constant-v0} ids
        work (they raised before).
  \item The environments no longer call \code{np.random.seed(42)} or
        \code{torch.manual\_seed(42)}; the random topology uses a local
        generator (\code{topology\_seed}) and is unchanged.
  \item The PyTorch backend accepts \code{device="auto"}, and its phase
        coherence and bonus rules match the NumPy backend.
  \item The dynamics are vectorised (20 times faster at $N = 50$, 65 times with
        RK4) and the renderer was rewritten; the former \code{"rgb\_array"} path
        failed with matplotlib 3.10 and later.
  \item \code{env\_lib.kos\_env.register()} is deprecated; the ids are
        registered when \pkg{env\_lib} is imported.
\end{itemize}

\endgroup
